\begin{quote}
Information associated with the components is stored in parameters and calculations. The engineers
set a parameter such as maximum power by providing a value and a unit. Provided calculations make use
of these parameters to e.g. let Virtual Satellite automatically calculate the average power consumption
of a component. By convention there are parameters with a special meaning such as the parameter
unit_quantity. This parameter denes how often the component is built into the spacecraft, whereas
mode_duration and duty_cycle represent the time a spacecraft remains in a specic system mode and
how much of that time a component is active. The parameter unit_active is important for calculating the
power consumptions in case of redundancies, when not all components are turned on at the same time.
[6] Visualizing the inputs and outputs of such calculations similar to UML block/interface diagrams was
successfully shown in previous work [14]. A language called ICML was proposed to model signals and
messages in the frame of interfaces [7].

The Open Concurrent Design Tool (OCDT) is developed by the European Space Agency (ESA). It is
developed for the use in CE and in ESA's Concurrent Design Facility (CDF) in particular. It provides a
CDM called Space Engineering Information Model (SEIM) including capabilities for structural decom-
position as well as storing information in parameters. A set of recommended parameters is dened under
the name Space Engineering Reference Data Library (SERDL). Calculations are done in external tools
such as Microsoft Excel. They can exchange information with the data-model by a generic REST API.
[9, 3]
Virtual Spacecraft Design (VSD) is another tool developed by ESA. It is designed for introducing
MBSE into later project phases. It provides capabilities for hierarchical breakdown and structuring of
the system, as well as adding components. On top of that, it oers conguration control mechanisms.
[4] Information is stored in predened objects such as state machines or interfaces. Designing functional
electrical architecture is well supported by VSD. It allows dening InterfaceEnds on components and
Interfaces connecting the InterfaceEnds. Individual information can be stored in engineering categories.
[13]
SysML provides functionality to model block-diagrams together with ports and interfaces. Blocks
within a block-diagram can be dened once and then reused and nested into other blocks, thus allowing
a hierarchical decomposition. The blocks can communicate to other blocks through ports. There are
dierent types of ports such as standard ports and ow ports. [12] Standard ports can either provide
or consume a service. Signals and data ows should be modeled using ow ports. There is no clear cut
denition when to use which one. [15]

Virtual Satellite 4 is the latest development. It is intended to support the whole lifecycle of a space-
craft. Since it is impossible to foresee all upcoming requirements on a data-model, it allows extending
the data-model along the lifecycle. The Generic System Engineering Language (GSEL) allows dening
new concepts for such a CDM extension. Concepts are activated when needed. A concept may contain
new structural elements for decomposition, new parameters to store information and calculations to
process them. The parameters are dened using an extended version of engineering categories. [5] These
categories are similar to the type-object and dynamic-template pattern [8, 11]. Engineers can dene a
category together with properties. These categories can then be instantiated on individual components
at run-time, where the engineers can now dene the value of the properties [13]. Based on an Eclipse
and EMF infrastructure, new concepts can be shipped with additional implemented functionality and
installed as new plug-ins [5].
\cite{review-1277}
\end{quote}

\begin{quote}
At CNES, PASO (Orbital Systems Architecture
Plateau) is leading space systems preliminary
phases and aims at formulating missions needs
and pre-sizing systems in order to conclude on
their feasibility. This process is conducted
accordingly to the principles of concurrent
engineering, thus involving all systems
architects and thematic experts from the very
beginning of the projects.
The objective of this paper is to describe the
approach through the prism of an electrical
architect, that contributes to needs analysis,
realization of power budgets, system sizing and
energetic simulations through a specific use
case. All this is achieved thanks to a precise
workflow involving multiple digital tools, with
continuity and interoperability made possible
between engineers and their tools by the
standardization of formats and files.

As an electrical architect specifically, the
objective of any project preliminary phase
would be:
- To define preliminary power budgets,
- To pre-size solar panels and batteries,
- To perform technological trade-offs,
- To simulate the on-board power usage.
Power simulations are a very good way to
compare technologies efficiently (solar cells
and battery cells for instance) and to refine the
power subsystem sizing. In order to be able to
perform energy balance simulations, several
elements are needed:
- A model of the power sub-system:
battery and solar generator assemblies,
power bus topology, type of
architecture, regulators, etc.
- All energy flows that leads the solar
generator to generate power: solar flux,
earth albedo flux, earth infrared flux.
- A power profile: the system power
consumption over time.
The model is defined in the Opalis software. It
is a proprietary but open tool developed by
CNES that allows performing power
simulations. The tool includes precise
modelling of electrical architectures including
batteries, solar panels, regulators and bus
topologies by taking into account the specific
features of each element such as losses,
temperature impacts, etc. Such parameters can
all be specified in the model tab of the tool.

The power profile is generated using IDM-CIC
by adding the IDM-Power plugin to it. This
plugin allows to allocate the power
consumptions described in the IDM-CIC power
budget to system modes. It then allows to obtain
a temporal view of the system’s power
consumption.

0) Preliminary work:
a. Defining every sub-system and
their constituting pieces of
equipment in IDM-CIC.
b. Defining the orbital scenario
using Simu-CIC.
1) Attributing a consumed/dissipated
power to each piece of equipment in
order to obtain the global power budget
of the system (including margins).
2) Converting the power budget into a
power profile (system power
consumption over time) using IDM-
Power plugin.
3) Importing power profile into the power
simulation software (Opalis).
4) Converting ephemeris into energy flows
(solar flux, earth infrared flux, earth
albedo flux) using the Flows Manager
plugin.
5) Defining the simulation model
Running energy balance simulations
using Opalis
\cite{review-0289}
\end{quote}

% Requirement
\begin{quote}
If left buried in the database of a model, the information is of little value. To be able to view or analyze the
information, reports are built.
\cite{review-1438}
\end{quote}

\begin{quote}
In many MBSE efforts, system models are built without a clear problem to be solved, and as a result, there are
cases where the models created are not usable or the MBSE effort itself is terminated, and the importance of "Problem
Framing" is suggested[16]. The proposed framework includes a step to carefully consider the issues to be solved
through the system model before building the system model, which is common with reference[16].
As a benefit of MBSE, a number of methods have been proposed to automatically generate SE deliverables from
the system model, such as "documents" required for specifications and review boards, and FTAs and FMEAs required
for risk analysis [14][15]. Compared to traditional document-based development, MBSE enhances the accumulation
and exchange of information necessary for systems engineering by structuring and centrally managing information
about the system.
\cite{handpicked-0004}
\end{quote}

% Potential PA candidate
\begin{quote}
Proposes an approach to CubeSat Mission Design that is both Concurrent and Model-based.
They state that such an approach has been adopted at a CubeSat Concurrent Design Facility, with roles including Project Manager, Systems Engineer, Mission Analyst, EPS Engineer, and a Risk and Cost Engineer.
The Project Manager is given an issue tracker software as a tool; the Systems Engineer has a Data Exchange System, a spreadsheet for sizing, budgeting, and change tracking; and the EPS Engineer has a database of commercial-off-the-shelf (COTS) components and their physical, functional, environmental, and power properties.
\cite{review-0529}
\end{quote}